\markboth{EVALUATION}{EVALUATION}
\section{Evaluation Framework}\label{sec:evaluation}

Evaluating a multi-agent conversational system requires assessing several dimensions at once: does it call the right tools, ground answers in fetched data, maintain a supportive tone, and let the user achieve their goal across multiple turns? This section describes the evaluation framework---conversation simulation, persona design, and scoring methodology.

\subsection{Simulation Pipeline}\label{sec:evaluation:framework}

We adopt MLflow's end-to-end conversation simulation framework~\citep{mlflow_multiturn_2025} for automated, reproducible evaluation of multi-turn agent systems. The four LLM-judge scorers rest on the finding of \cite{zheng_judging_2023} that a strong model grading open-ended dialogue agrees with human preferences about as often as humans agree with each other ($\sim$85\,\% GPT-4--human agreement, exceeding the 81\,\% human--human rate on their MT-Bench data), making LLM-as-a-judge a scalable proxy for human evaluation; their catalogue of judge biases (position, verbosity, self-enhancement) also motivates our model-separation safeguard (Section~\ref{sec:evaluation:models}). The scorer design further draws on \cite{shankar_validates_2024}, who showed that LLM-based evaluators must be aligned with human preferences to be reliable, and our deterministic grounding scorer follows the $\tau$-bench approach of \cite{yao_taubench_2024}, which evaluates tool-agent-user interaction by comparing system state against annotated goals rather than chat-text judgement alone---motivating our combination of LLM judges with a trace-based grounding scorer. The framework has three components: a \textbf{conversation simulator} that uses an LLM to role-play a user persona across a coherent multi-turn conversation; a set of \textbf{scorers} (LLM-based judges and deterministic code checkers) evaluating each conversation along defined quality dimensions; and an \textbf{experiment tracker} (MLflow) recording every conversation, trace, scorer verdict, and aggregate metric for reproducible comparison across runs. The pipeline is fully automated: a single command (\texttt{python -m evaluation.run\_e2e}) starts the simulator, drives all persona conversations against the live Training Copilot, scores them, and generates an HTML report.

\subsection{Persona Design}\label{sec:evaluation:personas}

We define ten synthetic personas across two user archetypes, directly informed by the market analysis in Section~\ref{sec:market:users}, designed to exercise both user sophistication levels and all five specialist domains. \textbf{Type A: ambitious triathletes} (5 personas)---structured, metrics-driven endurance athletes who train by plan, wear Garmin watches, and expect precise, data-driven decisions---each pursue a different multi-turn goal: Julian (recovery readiness for a key brick session), Priya (training-load trend analysis and taper planning), Marco (post-workout execution review with HR zones and splits), Lena (calendar-and-weather scheduling, with calendar writes), and Tom (evidence-based exercise-science guidance with citations). \textbf{Type B: hobby cyclists} (5 personas)---recreational riders who train by feel and prefer plain-language guidance---comprise Sophie (scenic loop with a caf\'e stop), Ben (weather go/no-go for a casual ride), Mara (trail discovery with elevation context), Carlos (year-over-year progress in plain language), and Jana (group-ride loop planning). Table~\ref{tab:personas} in Appendix~\ref{sec:appendix:personas} lists all ten personas with their type, goal, and expected focus area.

Each persona combines four elements consumed by the \texttt{ConversationSimulator}: a \textit{goal}, a \textit{persona identity} (background, fitness level, communication style), \textit{simulation guidelines} (e.g.\ ``push for concrete specifics'', ``stay in character''), and \textit{expectations} (the expected focus area, used for scorer calibration). Common guidelines instruct the simulator to have a natural multi-turn conversation, react to the assistant's actual answer before moving on, and end once the goal is genuinely met or clearly unreachable. The simulator LLM (GPT-5.4-mini) role-plays each persona for up to five turns.

\subsection{Scoring Dimensions}\label{sec:evaluation:scorers}

Each conversation is scored along five dimensions, combining LLM-based judges with a deterministic code scorer (Table~\ref{tab:scorers}).

\begin{table}[ht]
\begin{center}
\begin{minipage}{\textwidth}
\caption[Evaluation scoring dimensions]{Evaluation scoring dimensions. Four LLM judges run on GPT-5.4-nano; one deterministic scorer inspects tool-call traces directly.}\label{tab:scorers}
\begin{tabularx}{\textwidth}{p{3.5cm}lX}
\toprule
\textbf{Scorer} & \textbf{Type} & \textbf{Assessment} \\
\midrule
Conversation Completeness & LLM judge & Were the user's questions fully and satisfactorily answered? \\
\addlinespace
User Frustration & LLM judge & Did the user become frustrated, and did the agent worsen or mitigate it? \\
\addlinespace
Safety & LLM judge & Is the content safe and free of harmful advice? \\
\addlinespace
Supportive Coaching Tone & LLM judge & Does the assistant maintain a supportive, encouraging tone throughout, even when the user pushes back? \\
\addlinespace
Grounded in Real Data & Deterministic & Did the Copilot actually use its MCP tools to fetch real data, or did it answer from parametric knowledge alone? \\
\botrule
\end{tabularx}
\end{minipage}
\end{center}
\end{table}


The fifth scorer, \texttt{grounded\_in\_real\_data}, inspects the tool-call \textit{spans} reconstructed in each turn's MLflow trace rather than asking an LLM to guess whether tools were used---it counts tool calls, identifies which were invoked and whether they succeeded, and produces a per-turn breakdown. The verdict is binary (``yes'' if at least one tool was called across the conversation), but the rationale gives granular visibility into grounding quality; a factual ``did it call tools or not?'' question is far more reliably answered from the trace than guessed from chat text. The \texttt{supportive\_coaching\_tone} scorer implements a \texttt{ConversationalGuidelines} assertion that the assistant remain supportive even when the user is impatient, sceptical, or pushing for specifics, never dismissive or robotic---testing resilience under adversarial behaviour that \cite{barger_artificial_2025} identified as critical for AI coaching acceptance.

One capability these five scorers do not yet isolate is \textit{citation faithfulness}: whether the fitness specialist's evidence-based claims are actually supported by the passages it retrieved, rather than fluently invented. The Tom persona exercises this qualitatively (he demands and scrutinises sources), but a dedicated deterministic scorer---checking that every citation the answer surfaces corresponds to a passage present in that turn's RAG-tool span, analogous to how \texttt{grounded\_in\_real\_data} reads MCP spans---is a planned sixth dimension we report as future work. Complementing the automated scorers, a planned one-click \textit{report button} on each chat answer (Section~\ref{sec:agents:observability}) will let users flag responses against the same MLflow run ID; these human labels give a validation set against which the LLM judges themselves can be checked for the alignment \cite{shankar_validates_2024} argue is prerequisite to trusting them.

\subsection{Trace Reconstruction}\label{sec:evaluation:trace}

A technical challenge in evaluating the multi-agent system is that the deep tool-call spans produced by the agent server processes live in a separate MLflow experiment from the evaluation's conversation traces. The evaluation process addresses this by reconstructing the Copilot's specialist and tool-call structure from the \texttt{trace} dict returned by \texttt{orchestrator.run()} and emitting them as real child spans within the evaluation trace (Figure~\ref{fig:trace_reconstruction}).

\begin{figure}[ht]
\centering
\resizebox{\textwidth}{!}{%
\begin{tikzpicture}[
    node distance=0.3cm and 0.5cm,
    every node/.style={font=\small},
    spanbox/.style={draw, thick, rounded corners=2pt, minimum height=0.55cm, align=center, font=\small},
    rootspan/.style={spanbox, draw=tcblue, fill=tcblue!8, minimum width=9cm},
    agentspan/.style={spanbox, draw=tcyellow!80!black, fill=tcyellow!8, minimum width=4cm},
    toolspan/.style={spanbox, draw=tccyan, fill=tccyan!6, minimum width=3.2cm},
    arr/.style={-{Stealth[length=4pt]}, thick, tcgrey!50},
]
    \node[rootspan] (root) {\texttt{fitdash\_copilot} (AGENT span)};

    \node[agentspan, below left=0.6cm and 0.3cm of root] (ag1) {\texttt{agent:recovery}};
    \node[agentspan, below right=0.6cm and 0.3cm of root] (ag2) {\texttt{agent:context}};

    \node[toolspan, below left=0.5cm and -0.2cm of ag1] (t1) {\texttt{garmin\_\_sleep}};
    \node[toolspan, right=0.15cm of t1] (t2) {\texttt{garmin\_\_hrv}};

    \node[toolspan, below left=0.5cm and -0.2cm of ag2] (t3) {\texttt{weather\_\_forecast}};
    \node[toolspan, right=0.15cm of t3] (t4) {\texttt{calendar\_\_events}};

    \draw[arr] (root.south) -- (ag1.north);
    \draw[arr] (root.south) -- (ag2.north);
    \draw[arr] (ag1.south) -- (t1.north);
    \draw[arr] (ag1.south) -- (t2.north);
    \draw[arr] (ag2.south) -- (t3.north);
    \draw[arr] (ag2.south) -- (t4.north);

\end{tikzpicture}%
}% end resizebox
\caption[Reconstructed span tree in the evaluation trace]{Reconstructed span tree in the evaluation trace. Agent and tool spans are rebuilt from the Copilot's \texttt{trace} dict and nested under the root evaluation span. Each tool span carries four custom attributes: \texttt{fitdash.tool} (name), \texttt{fitdash.tool\_ok} (success), \texttt{fitdash.tool\_error} (message), and \texttt{fitdash.duration\_ms} (latency).}\label{fig:trace_reconstruction}
\end{figure}


Each reconstructed tool span carries four custom attributes---\texttt{fitdash.tool} (the namespaced tool name), \texttt{fitdash.tool\_ok} (boolean success), \texttt{fitdash.tool\_error} (error message if any), and \texttt{fitdash.duration\_ms} (tool-call latency)---and the grounding scorer reads only these, never the chat text, making its verdict deterministic and model-independent.

\subsection{Model Separation}\label{sec:evaluation:models}

A critical design decision is separating the models used for evaluation from those used by the system under test: the evaluation framework uses official OpenAI models (GPT-5.4-mini for persona simulation and report generation; GPT-5.4-nano for the four LLM judges), while Training Copilot itself runs on the KIT university gateway with its own model configuration. The \texttt{apply\_openai\_routing()} function rewrites the evaluation process's environment to use the official OpenAI key and base URL, ensuring judges and system under test share no weights, biases, or provider-specific behaviour.

\subsection{Metrics and Reporting}\label{sec:evaluation:metrics}

Each LLM judge returns a per-conversation verdict, which we aggregate as a \textit{pass rate}---the fraction of the ten conversations the judge marks acceptable---reported overall and split by persona type; the deterministic \texttt{grounded\_in\_real\_data} scorer aggregates the same way over its binary per-conversation verdict. Alongside the five quality scorers we log two operational measures per conversation---the number of turns taken to reach the goal and end-to-end latency---so conversational quality can be read against cost. MLflow rolls these into experiment-level metrics, and the generated HTML report renders them next to every transcript and trace for inspection. Table~\ref{tab:results} gives the reporting structure; its cells are placeholders pending the final run. We stress that ten personas of at most five turns is a \textit{formative} sample: it is sized to exercise every specialist domain and surface failure modes, not to establish statistical significance, and we report it as a characterisation of system behaviour rather than a confirmatory benchmark.

\begin{table}[ht]
\begin{center}
\begin{minipage}{\textwidth}
\caption[Aggregate evaluation results (placeholder)]{Aggregate evaluation results across the ten persona conversations, reported as the mean scorer pass rate (fraction of conversations the scorer marks acceptable; for \texttt{grounded\_in\_real\_data} the fraction that issued at least one successful tool call). Values are placeholders (\textbf{TBD}) to be populated from the final run before submission.}\label{tab:results}
\begin{tabularx}{\textwidth}{Xllc}
\toprule
\textbf{Scoring Dimension} & \textbf{Type} & \textbf{Triathletes} & \textbf{Overall} \\
 & & \textbf{/ Cyclists} & \textbf{(n=10)} \\
\midrule
Conversation Completeness & LLM judge & TBD / TBD & TBD \\
\addlinespace
User Frustration (resilience) & LLM judge & TBD / TBD & TBD \\
\addlinespace
Safety & LLM judge & TBD / TBD & TBD \\
\addlinespace
Supportive Coaching Tone & LLM judge & TBD / TBD & TBD \\
\addlinespace
Grounded in Real Data & Deterministic & TBD / TBD & TBD \\
\midrule
Mean turns to goal & --- & TBD / TBD & TBD \\
Mean end-to-end latency (s) & --- & TBD / TBD & TBD \\
\botrule
\end{tabularx}
\end{minipage}
\end{center}
\end{table}

\FloatBarrier

\subsection{Planned Ablation: Multi-Agent vs.\ Single-Agent}\label{sec:evaluation:ablation}

The design's central hypothesis---that scoping specialists to narrow tool sets improves tool selection over one agent holding all $\sim$45 tools---was observed only informally during early prototyping (Section~\ref{sec:architecture:rationale}). To substantiate it, the same ten persona conversations will be run against a single-agent baseline that exposes every server's tools to one ReAct agent, holding the model, MCP servers, and personas fixed so that only the agent topology varies. We compare correct-tool-selection rate (judged against each persona's expected focus area), conversation completeness, grounding, and mean latency (Table~\ref{tab:ablation}). This directly tests whether the coordination overhead the multi-agent design adds (Section~\ref{sec:discussion:tradeoffs}) is repaid in answer quality.

\begin{table}[ht]
\begin{center}
\begin{minipage}{\textwidth}
\caption[Multi-agent vs.\ single-agent ablation (placeholder)]{Planned ablation isolating the effect of domain-scoped multi-agent decomposition. Both configurations answer the same ten persona conversations with the same models and MCP servers; only the agent topology differs. The single-agent baseline exposes all servers' tools to one ReAct agent. Values are placeholders (\textbf{TBD}) to be populated from the final run.}\label{tab:ablation}
\begin{tabularx}{\textwidth}{Xcccc}
\toprule
\textbf{Configuration} & \textbf{Correct tool} & \textbf{Complete-} & \textbf{Grounded} & \textbf{Mean} \\
 & \textbf{selection} & \textbf{ness} & & \textbf{latency (s)} \\
\midrule
Single agent (all $\sim$50 tools, one scope) & TBD & TBD & TBD & TBD \\
\addlinespace
Multi-agent (5 domain-scoped specialists) & TBD & TBD & TBD & TBD \\
\botrule
\end{tabularx}
\end{minipage}
\end{center}
\end{table}

\FloatBarrier

\subsection{Expected Outcomes}\label{sec:evaluation:expected}

We expect the evaluation to confirm high grounding scores, since the orchestrator has no MCP access of its own and specialists must call tools to access data; specialist coverage, since the ten personas collectively exercise all five domains; coaching-tone resilience, since the ambitious-triathlete personas are deliberately demanding (Julian is ``slightly impatient with hand-waving''; Tom ``distrusts influencer fitness tips''); and, in the ablation, a higher correct-tool-selection rate for the scoped multi-agent configuration than for the single-agent baseline. We anticipate two weaknesses: multi-step sequential queries, where a later specialist's input depends on an earlier one's output, may show higher latency and occasional context loss, and the fitness specialist's public-domain corpus may not satisfy users expecting contemporary sports-science references. Quantitative results will be reported in the final submission and will inform iterative improvements to the prompt architecture and specialist scoping.
